% Paste into main.tex — Results subsection (6E7B supplementary MD)
% Source: analysis/summary.md (2026-06-15, 3 × 200 ns replicates)

\subsection{CPPF binding in the 6E7B ($\beta$-GMPCPP) conformational state}

To address whether CPPF binding is maintained in the microtubule-lattice-related $\beta$-tubulin conformation (PDB 6E7B), we performed three independent 200~ns MD replicates of the CPPF--$\alpha\beta$-tubulin complex using the same cofactor-free protocol as the main-text 5IJ0 simulations.
Across all replicates, CPPF remained stably associated with the heterodimer: over the final 50~ns (150--200~ns), the backbone RMSD was $0.295 \pm 0.026$~nm (mean $\pm$ s.d.\ across replicates; per-replicate values 0.270, 0.295, and 0.321~nm) and the minimum CPPF--protein heavy-atom distance was $0.203 \pm 0.014$~nm (per-replicate means 0.209, 0.214, and 0.188~nm).
These distances are comparable to the main-text 5IJ0 dimer ensemble ($\sim$0.19~nm; range 0.14--0.24~nm), indicating that CPPF retains a stable binding mode in the alternative $\beta$-tubulin conformational state.
Free-energy landscape analysis placed the global minimum at RMSD 0.27~nm and $R_g$ 2.99~nm, consistent with a single well-defined bound basin (Supplementary Fig.~X).
Trajectory files are available at \url{https://huggingface.co/datasets/HUB_NAMESPACE/MD-trajectories-CPPF-tubulin-heterodimer-and-monomers}.
